\section{Conclusions and Future Work}

We presented SimPLe, a model-based reinforcement learning approach that operates directly on raw pixel observations and learns effective policies to play games in the Atari Learning Environment. Our experiments demonstrate that SimPLe learns to play many of the games with just $100$K interactions with the environment, corresponding to 2 hours of play time. In many cases, the number of samples required for prior methods to learn to reach the same reward value is several times larger.

Our predictive model has stochastic latent variables so it can be applied in highly stochastic environments. Studying such environments is an exciting direction for future work, as is the study of other ways in which the predictive neural network model could be used. Our approach uses the model as a learned simulator and directly applies model-free policy learning to acquire the policy. However, we could use the model for planning. Also, since our model is differentiable, the additional information contained in its gradients could be incorporated into the reinforcement learning process. Finally, the representation learned by the predictive model is likely be more meaningful by itself than the raw pixel observations from the environment. Incorporating this representation into the policy could further accelerate and improve the reinforcement learning process.

While SimPLe is able to learn more quickly than model-free methods, it does have limitations. First, the final scores are on the whole lower than the best state-of-the-art model-free methods. This can be improved with better dynamics models and, while generally common with model-based RL algorithms, suggests an important direction for future work. Another, less obvious limitation is that the performance of our method generally varied substantially between different runs on the same game. The complex interactions between the model, policy, and data collection were likely responsible for this. In future work, models that capture uncertainty via Bayesian parameter posteriors or ensembles \citep{trpo_ensemble, Chua18} may improve robustness.
Finally, the computational and time requirement of training inside world model are substantial (see Appendix \ref{sec:architecture-details}), which makes developing lighter models an important research direction.

In this paper our focus was to demonstrate the capability and generality of SimPLe only across a suite of Atari games, however, we believe similar methods can be applied to other environments and tasks which is one of our main directions for future work. As a long-term challenge, we believe that model-based reinforcement learning based on stochastic predictive models represents a promising and highly efficient alternative to model-free RL. Applications of such approaches to both high-fidelity simulated environments and real-world data represent an exciting direction for future work that can enable highly efficient learning of behaviors from raw sensory inputs in domains such as robotics and autonomous driving.

\subsection*{Acknowledgments} We thank Marc Bellemare and Pablo Castro for their help with Rainbow and Dopamine. The work of Konrad Czechowski, Piotr Kozakowski and Piotr Miłoś was supported by the Polish National Science Center grants UMO-2017/26/E/ST6/00622. The work of Henryk Michalewski was supported by the Polish National Science Center grant UMO-2018/29/B/ST6/02959. This research was supported by the PL-Grid Infrastructure. In particular, Konrad Czechowski, Piotr Kozakowski, Henryk Michalewski, Piotr Miłoś and Błażej Osiński extensively used the Prometheus supercomputer, located in the Academic Computer Center
Cyfronet in the AGH University of Science and Technology in Kraków, Poland.
Some of the experiments were managed using \url{https://neptune.ai}. We would like to thank the Neptune team for providing us access to the team version and technical support.